\section{Le Chat：快推理、工具层与产品反馈飞轮}

消费者产品在访谈中承担三种角色：展示模型和工具能力，产生真实任务分布，收集失败反馈。课堂同时指出，极致 inference optimization 可能降低模型更新频率或架构兼容性。这是典型系统 trade-off：速度并非免费，它可能要求固定 runtime、专用 kernel、受限模型族或更复杂的发布流程。

\subsection{产品为什么是训练系统的一部分}

用户请求、tool trace 和显式评分可以揭示 benchmark 看不到的问题：哪些查询需要 web search，哪些代码任务常失败，哪些语言或行业表现不稳定。产品团队随后把这些失败变成 evaluation slice 和训练优先级，形成 data flywheel。

\lecturefigure{12-product-feedback-flywheel.png}{Le Chat 类产品把 serving、工具与训练优先级连接成反馈飞轮。}{本地字幕 24:45--27:40；概念重绘。}

\begin{knowledgebox}{读图：反馈必须先变成 failure taxonomy}
简单 thumbs up/down 含义模糊：低分可能来自事实错误、语气、工具超时、界面延迟或用户目标变化。可靠飞轮要保留 request context、模型与工具版本、延迟、最终结果和人工分类，再决定是改数据、模型、prompt、tool 还是 serving。
\end{knowledgebox}

\begin{warningbox}{“收集产品数据”必须有治理边界}
产品日志可能包含个人信息、企业机密或受监管数据。训练回流需要 purpose limitation、retention policy、访问控制、用户选择、去标识化和数据 lineage。反馈飞轮不是把所有对话自动加入训练集。
\end{warningbox}

\teachervoice{课堂还给出一个具体产品信号：某一时点英语请求中代码任务占比很高，于是团队会增加 code generation 投入。这个例子说明，roadmap 不只来自领导判断，也可以来自可观测的任务分布；但观测数据仍要校正用户群偏差。}

\subsection{本章小结}

产品既是服务界面，也是任务采样器和失败观测站。高质量反馈飞轮必须把用户信号分解为可定位问题，并在隐私治理下回到评测和训练。

